\subsection{Increasing clocks and based sphere parameters}
The compact resolution also gives increasing changes of a scalar
parameter and relative regularizations on a sphere.  The first
construction keeps a prescribed finite set of times fixed; the second
keeps the basepoint fixed.


Over the reals, the Hausdorff-resolution theorem of
Kocel-Cynk--Paw\l ucki--Valette provides a Lipschitz parameter
resolution for a bounded family of compact sets
\cite[Main Theorem, author manuscript p.~1]{KocelCynkPawluckiValette2014}.
Valette's generic Lipschitz proposition instead gives joint bounds
on compact subsets of finitely many parameter pieces
\cite[Proposition~5.7, pp.~25--26; scope on p.~22]{Valette2022RegularVectors}.
The scalar consequence below gives one increasing clock fixing
any prescribed finite set. It follows over every real closed field
from the finite cube resolution proved above. The sphere and based
homeomorphism-family consequences below also hold over every such field.

\begin{corollary}[An increasing definable clock]
\label{cor:increasing-definable-clock}
Let $M$ be an o-minimal expansion of a real closed field, let
$Z\subset M^n$ be closed, bounded and definable, and let
$F:[0,1]\times Z\to M$ be continuous and definable.  Suppose
\[
 |F(t,z)-F(t,z')|\le L\|z-z'\|_2
 \qquad(t\in[0,1],\ z,z'\in Z)
\]
for some $L\in M_{\ge0}$.  For every finite set $E\subset[0,1]$
there is an increasing definable Lipschitz homeomorphism
$h:[0,1]\to[0,1]$, fixing $E$, and $K\in M_{\ge0}$ such that
\[
 |F(h(s),z)-F(h(s'),z')|
 \le K|s-s'|+L\|z-z'\|_2.
\]
If $(A_t)_{t\in[0,1]}$ is a Hausdorff-continuous definable family
of nonempty closed bounded sets contained in one bounded box, such
an $h$ also satisfies
\[
 d_H(A_{h(s)},A_{h(s')})\le K|s-s'|.
\]
The constants $L,K$ may be arbitrary elements of $M$.
\end{corollary}

\begin{proof}
We first extract an increasing clock from the finite cube resolution.
Suppose $\phi:K\to[0,1]$ is a continuous definable surjection from
a finite disjoint union of closed cubes, and that a metric-valued
family $G_t$ has a Lipschitz pullback $G_\phi$.  In each cube on
which $\phi$ is nonconstant, choose a point of minimum and a point
of maximum, and join them by an affine segment $\alpha:[0,1]\to K$.
The continuous definable function $\phi\alpha$ has image the full
interval between that cube's minimum and maximum.  Partition its
domain into finitely many closed intervals on each of which it is
constant or strictly monotone.

The nondegenerate image intervals of these monotone pieces cover
$[0,1]$.  Indeed, their union is closed.  Every value omitted from
the union belongs to the finite set consisting of constant-piece
values and images of constant cubes.  A nonempty relatively open
subset of $[0,1]$ cannot be a finite set, so nothing is omitted.
This argument includes zero-dimensional cubes.

Partition $[0,1]$ at all endpoints of these image intervals and all
points of $E\cup\{0,1\}$.  Each successive interval $[a_i,a_{i+1}]$
lies in the image of one strict monotonicity piece: take an image
interval containing its midpoint, and use the fact that every
endpoint of every image interval is a partition point.  Let
$u_i,v_i$ be the preimages of $a_i,a_{i+1}$ in that piece, and put
\[
 h(s)=\phi\alpha_i\left(
 u_i+\frac{v_i-u_i}{a_{i+1}-a_i}(s-a_i)\right)
 \qquad(a_i\le s\le a_{i+1}).
\]
If the piece is decreasing, then $u_i>v_i$, so the displayed
composition still increases.  Its endpoint values are $a_i$ and
$a_{i+1}$.  The formulas therefore join to a continuous strictly
increasing definable bijection of $[0,1]$, hence a homeomorphism,
fixing every partition point.

The argument of $\phi$ on each such interval is an affine path in
one cube with a finite field-valued Lipschitz constant.  The
Lipschitz bound for $G_\phi$ gives a bound for $G_h$ on that interval.
Take the maximum of the finitely many bounds.  For parameters in
different intervals, split the interval between them at the
partition points and apply the triangle inequality.  The lengths
of the pieces add to the original parameter distance, so this
maximum is a global bound.  The same reasoning shows that $h$ is
Lipschitz whenever $\phi$ is Lipschitz.  Only finite monotonicity,
definable extrema and the definable intermediate value theorem
have been used; all are valid over $M$.

If $Z$ is empty, the function assertion is vacuous and $h$ may be
the identity.  Otherwise place $Z$ in a closed bounded coordinate
box $Q$, and extend the function to this box by
\[
 \widetilde F(t,x)=
 \min_{z\in Z}\bigl(F(t,z)+L\|x-z\|_2\bigr).
\]
Definable compactness gives the minimum.  The given spatial
inequality proves equality with $F$ on $Z$, and the triangle
inequality gives the same spatial constant $L$ on $Q$.  Moreover,
\[
 |\widetilde F(t,x)-\widetilde F(t',x')|
 \le \sup_{z\in Z}|F(t,z)-F(t',z)|+L\|x-x'\|_2.
\]
Definable uniform continuity on the closed bounded domain proves
joint continuity.  Normalize the box and the bounded range by
affine changes of coordinates and apply Theorem~\ref{result-18}.
Norm equivalence changes only the bound used for that theorem.
It gives a finite cube resolution for which the pullback is
jointly Lipschitz.

Apply the preceding clock construction to $G_t=F(t,-)$ with the
supremum metric.  The resulting parameter estimate is uniform
in $z$.  Combining it with the original spatial estimate gives
the asserted inequality with the original coefficient $L$.
For the compact-set assertion, apply the Hausdorff consequence of
Theorem~\ref{result-18} and the same construction with $G_t=A_t$.
No point matching between different compact sets is required.
\end{proof}

For an explicit example with a fixed intermediate time, let
$F(t,z)=\sqrt t+Lz$ on $[0,1]^2$ and require $h(1/2)=1/2$.
Writing $a=1/\sqrt2$, the clock
\[
 h(s)=
 \begin{cases}
  2s^2,&0\le s\le1/2,\\
  \bigl(a+2(1-a)(s-1/2)\bigr)^2,&1/2\le s\le1
 \end{cases}
\]
is increasing and fixes $0,1/2,1$.  The function $\sqrt{h(s)}$
is affine on each half, with slopes $\sqrt2$ and $2-\sqrt2$.
Thus the parameter constant may be taken to be $\sqrt2$, while
the spatial constant stays $L$.

\begin{corollary}[Hausdorff regularization on a based sphere]
\label{cor:based-sphere-hausdorff}
Let $R$ be a real closed field with a fixed o-minimal expansion.
In this statement compact means closed, bounded and definable, and
all norms are maximum norms.  Let
$\Sigma^k=\partial[-1,1]^{k+1}$, $k\ge1$, with a chosen basepoint $b$,
and let $(C_p)_{p\in\Sigma^k}$ be a Hausdorff-continuous definable
family of nonempty compact sets with bounded incidence.
There is a definable Lipschitz map $g:\Sigma^k\to\Sigma^k$ which
fixes $b$, is definably homotopic to the identity through maps fixing
$b$, and makes $p\mapsto C_{g(p)}$ Hausdorff-Lipschitz.
The Lipschitz constants may be arbitrary elements of $R$.

On a parameter cube, a Hausdorff-Lipschitz family on the entire
boundary admits the corresponding regularization fixing that
whole boundary.  In particular, let $F$ be a closed bounded
definable set and let $(U_p)$ be a jointly continuous definable
family of homeomorphisms of $F$ on a parameter cube, with
$U_p=\operatorname{id}_F$ on its entire boundary.
There is a definable homotopy relative to that boundary through
homeomorphism families to one with Hausdorff-Lipschitz graphs.
Every spatial condition common to the original family persists.
The full map and its inverse family are jointly continuous in
the field topology throughout the homotopy.  The empty $F$ case
is allowed separately.
\end{corollary}

\begin{proof}
The cube assertion is the relative conclusion of
Theorem~\ref{result-18}.  We derive the sphere assertion using
finite cells and estimates over $R$.

First, for $d\ge2$, any $x,y\in\partial[-1,1]^d$ can be joined
in that boundary by at most three segments, each lying in one
facet, with total maximum-norm length at most
$3\|x-y\|_\infty$.  Choose coordinates $x_i=\epsilon$ and
$y_j=\delta$ with $\epsilon,\delta\in\{-1,1\}$.
If $i\ne j$, change the $j$-coordinate of $x$ to $y_j$ and
then join to $y$.  Each of the two segments has length at most
$\|x-y\|_\infty$ and lies in one of the specified facets.
If $i=j$ with the same sign, use the common-facet segment.
If the signs are opposite, the distance is $2$; choose a
different coordinate, change it to $1$ at $x$ and at $y$,
and use the resulting three segments, each of length at most $2$.

Consequently a map into any $R$-valued metric space that is
$L$-Lipschitz on each facet and agrees on their intersections
is $3L$-Lipschitz on the whole boundary.  Affine normalization
gives the same conclusion for a rectangular cell, with a
constant depending on its positive side lengths.
For a one-dimensional cell boundary, the two values are
Lipschitz directly: divide their distance by the positive
distance between the endpoints.

We also need a Lipschitz extension of a boundary map.
For an $L$-Lipschitz map
$a:\partial[-1,1]^d\to[-1,1]^d$, put
\[
 E(0)=0,\qquad
 E(x)=\|x\|_\infty\,a(x/\|x\|_\infty)\quad(x\ne0).
\]
It extends $a$ and takes the cube into itself.
For $r=\|x\|_\infty\ge s=\|y\|_\infty>0$ and
$D=\|x-y\|_\infty$,
\[
 s\|x/r-y/s\|_\infty
 \le (s/r)D+(s/r)(r-s)\le2D.
\]
Since $r-s\le D$ and $\|a\|_\infty\le1$, this proves
$\|E(x)-E(y)\|_\infty\le(1+2L)D$.
The case $s=0$ is immediate.  Affine conjugation gives
the corresponding extension on every nondegenerate
rectangular cell.

In ambient coordinate $i$, subdivide $[-1,1]$ at the distinct
elements of $\{-1,b_i,1\}$.  The cells lying in $\Sigma^k$
form a finite rectangular complex, and $b$ is a vertex.
Construct $g$ in increasing cell dimension, fixing every
vertex.  Require that $g$ take every cell into itself, agree
on common faces, and that both $g$ and the pulled family be
Lipschitz on each constructed cell.

Suppose the construction is complete on the proper faces
of a cell $Q$.  The pasted map on $\partial Q$ and its
pulled family are Lipschitz by the boundary estimate above,
with the separate two-point argument when $\dim Q=1$.
Cone the boundary map to a definable Lipschitz
$E_Q:Q\to Q$.  The family
\[
 x\longmapsto C_{E_Q(x)}
\]
has nonempty compact fibres, bounded incidence, and
Hausdorff continuity, and its entire boundary family is
Hausdorff-Lipschitz.  Apply the relative cube theorem in
affine coordinates on $Q$.  It supplies a definable
Lipschitz self-map $h_Q$ fixing $\partial Q$ such that
$x\mapsto C_{E_Q(h_Q(x))}$ is Hausdorff-Lipschitz.
Set $g|_Q=E_Qh_Q$.  This agrees exactly with all earlier
faces and preserves $Q$.  Finite induction completes $g$.

These cellwise bounds are global.  On an original facet
of $[-1,1]^{k+1}$, split a straight segment at its finitely
many grid crossings.  Each piece lies in one closed
cell, and its lengths sum to the original distance.
The maximum of the finitely many cell constants therefore
bounds the restriction to the entire facet.  The preceding
boundary-path estimate in dimension $k+1\ge2$ gives the
global bounds on $\Sigma^k$, for both $g$ and the
Hausdorff metric of its pulled family.
All these constants lie in $R$; infinitesimal grid widths
only change the finite affine-normalization constants.

The formulas agree on closed cell intersections, so
finite closed pasting gives continuity.  For $p$ in
any cell $Q$, both $p$ and $g(p)$ belong to $Q$.  Hence
\[
 H(s,p)=(1-s)p+s\,g(p)
\]
stays in the same convex cell and thus in $\Sigma^k$.
It is a jointly continuous definable homotopy from the
identity to $g$, fixing all vertices and in particular $b$.

For the homeomorphism-family assertion, assume first
$F\ne\varnothing$ and write $P$ for its parameter cube.
Definable uniform continuity on the closed bounded
domain $P\times F$ gives Hausdorff continuity of
$\Gamma_p=\{(x,U_p(x)):x\in F\}$, since
\[
 d_H(\Gamma_p,\Gamma_q)
 \le \max_{x\in F}\|U_p(x)-U_q(x)\|_\infty.
\]
These graphs are nonempty closed bounded definable sets
with bounded incidence.  On $\partial P$ they are the
constant diagonal.  Apply the relative cube theorem to
obtain a boundary-fixed Lipschitz $h:P\to P$ for which
$\Gamma_{h(p)}$ is Hausdorff-Lipschitz, and put
\[
 U^s_p(x)=U_{(1-s)p+s h(p)}(x).
\]
Every member is one of the original homeomorphisms,
so every common spatial condition persists, and the
homotopy is relative to the whole parameter boundary.

For completeness, inverse continuity over $R$ follows
from definable extrema, without ordinary compactness.
If $\Psi:X\to Y$ is a continuous definable bijection and
$X$ is closed and bounded, then, for each $\epsilon>0$,
the nonempty set of pairs $(x,x')\in X^2$ with
$\|x-x'\|_\infty\ge\epsilon$ is closed bounded definable.
On it the continuous function
$\|\Psi(x)-\Psi(x')\|_\infty$ has a positive minimum:
the minimum is attained and injectivity excludes zero.
Thus $\Psi^{-1}$ is uniformly continuous; when the
set of pairs is empty the assertion is immediate.
The general-field extrema statement and uniform-continuity
fact are respectively Theorem~3.4 and Exercise~3.5 of
\cite[pp.~26--27; field convention pp.~9--10]{Coste1999}.

Apply this inverse argument to
\[
 (s,p,x)\longmapsto(s,p,U^s_p(x)),
\]
a continuous definable bijection of the closed bounded
definable set $[0,1]\times P\times F$ onto itself.
It gives the jointly continuous definable inverse
family, as required.  A zero-dimensional parameter
cube already gives a constant graph family.
For $F=\varnothing$ there is a unique family and
constant homotopy; its empty graphs have self-distance
zero when this convention is required.
\end{proof}

The parameter map in the sphere and cube assertions need not be
injective.  The conclusions control the graphs as compact sets;
they do not choose homeomorphisms matching their points.
